% Figure 4 manuscript-ready PRL-style text block

% Main-text paragraph
To determine whether finite-bias peaks alone can identify a chiral channel, we
constructed a non-chiral control kernel with symmetric peaks at $\pm V$,
tunable width, and tunable spectral weight, and evaluated it on the same
$\alpha$-$Z$-$\phi$ grid as the chiral branch. Although this control channel
reproduces finite-bias double peaks, it remains phase blind: the positive-bias
peak position shows no $\phi$-dependent drift, and the peak splitting shows no
coupling to $\alpha$. By contrast, the chiral channel exhibits both effects on
the same grid. In the present scan, the maximum drift of the positive-bias peak
position reaches $0.0928$, while the maximum $\alpha$-coupled change in the
peak splitting reaches $0.1856$; both quantities vanish in the control
channel. The two channels also differ in their barrier evolution. The chiral
branch is non-monotonic throughout the tested configuration set, whereas the
control channel retains a monotonic fraction of $0.67$. Finite-bias peaks are
therefore not sufficient by themselves. The discriminating feature is the
phase-sensitive locking among peak position, interface orientation, and barrier
evolution, which appears only in the chiral branch and is absent in the trivial
double-peak control.

% Figure caption
\caption{\textbf{Control discriminant for chiral finite-bias peaks.}
Comparison between the chiral finite-bias edge-state kernel and a non-chiral
control kernel with symmetric peaks at $\pm V$, tunable width, and tunable
spectral weight, evaluated on the same $\alpha$-$Z$-$\phi$ grid. The control
channel generates finite-bias double peaks but remains phase blind: the
positive-bias peak position does not drift with $\phi$, and the peak splitting
does not couple to $\alpha$. By contrast, the chiral channel exhibits both
phase-locked peak-position drift and $\alpha$-coupled peak splitting. In the
present scan, the maximum positive-peak drift is $0.0928$ for the chiral
channel and $0$ for the control, while the maximum $\alpha$-coupled change in
peak splitting is $0.1856$ for the chiral channel and $0$ for the control. The
chiral branch is also non-monotonic over the full tested configuration set,
whereas the control channel retains a monotonic fraction of $0.67$. Finite-bias
peaks alone are therefore insufficient: the manuscript claim rests on the
phase-sensitive coupling among peak position, interface orientation, and
barrier evolution, which the trivial double-peak control fails to reproduce.}
